\ssrno{РЕФЕРАТ}

Расчетно-пояснительная записка 64 с., 31 рис., 0 лист., 56 ист., 2 прил.

НЕЙРОННАЯ СЕТЬ, ОБНАРУЖЕНИЕ ПРИЗНАКОВ ГЕНЕРАЦИИ, ГЛУБОКОЕ ОБУЧЕНИЕ, КОМПЬЮТЕРНОЕ ЗРЕНИЕ, АНСАМБЛЬ МОДЕЛЕЙ, EFFICIENTNET, VISION TRANSFORMER, FACEFORENSICS.

Цель работы: разработка метода обнаружения признаков генерации на изображении с использованием нейронной сети.

В данной работе разрабатывается метод обнаружения признаков генерации на изображении с использованием нейронной сети. Проводится анализ предметной области и известных нейросетевых решений. Выполняется классификация и сравнительный анализ сверточных, капсульных нейронных сетей и трансформеров. Проектируется архитектура ансамбля нейронных сетей на основе EfficientNet и Vision Transformer. Реализуется программное обеспечение с графическим интерфейсом, и производится обучение модели на расширенном наборе данных FaceForensics++. Проводится исследование точности и устойчивости разработанного метода к сжатию изображений.

В результате работы разработан метод обнаружения признаков генерации на изображении, демонстрирующий точность 98,75\% на наборе данных FaceForensics и превосходящий рассмотренные аналоги по устойчивости к шуму. Создано приложение, позволяющее загружать изображения и получать вероятностную оценку наличия признаков генерации. Выполнено тестирование программного обеспечения и подтверждена его работоспособность на реальных сценариях.

